\subsection{Putting the results together}

Looking at the ten results side by side, one message stands out: \textbf{the
rules of the quiz shape student behaviour more than the time they spend.} Two
rules explain a lot of what we saw.

\begin{itemize}
  \item \textbf{The pass mark (H6).} Students keep trying until they reach
        4/10 and then stop. This is why almost every re-attempt path goes up,
        for weak and strong students alike (H3), and why so many students end
        at exactly 4.
  \item \textbf{The attempt delay (H7).} While re-attempts were allowed
        straight away, students could copy the answers they had just seen and
        get full marks in about a minute. These fast, perfect re-attempts are
        what flips the time--score relationship in H1, and they made ``very
        fast'' attempts look good when all attempts were mixed together.
\end{itemize}

A second message is that \textbf{speed means different things for different
students}. Strong students are not faster on hard questions (H4). It is the
weaker students who are fastest, and they are also the ones who leave questions
blank (H10) and finish in under a minute with low marks (H5). This matches the
idea of rapid, low-effort responding \citep{wise2005response} more than the
idea of slow, struggling students. It also fits the finding of
\citet{goldhammer2014time} that the meaning of time on task depends on the
person and the task.

The most stable thing in the data is \textbf{question difficulty} (H2). The
same questions are hard for every tier and in almost every half-year. The few
clear changes (for example in mid-2024 in Quizzes 2 and 3) are sudden, which
suggests the questions themselves were edited (H9).

Finally, \textbf{a single quiz is a weak signal of a student's level} (H8).
First-attempt grades in different quizzes are only weakly related
($\rho \approx 0.2$). Part of this is because the quizzes cover different
material and have very different difficulty, but it is also a reminder that
five two-mark questions give a noisy measure.

\subsection{What a teacher could do with these results}

\begin{enumerate}
  \item \textbf{Keep the enforced delay between attempts.} H7 shows it removed
        most of the ``copy the feedback'' behaviour. Showing the correct
        answers only after the quiz closes, or drawing each re-attempt from a
        pool of question variants, would make re-attempts even more useful.
  \item \textbf{Review Q1.4, Q1.5 and Q3.3.} Only 9--31\% of students get
        these right on the first try. Q1.5 is a special case: even High-tier
        students get it right only 16\% of the time, against 9\% for Low-tier
        students (Figure~\ref{fig:h2}b). A question that strong and weak
        students fail almost equally often does not separate them, and Moodle's
        own guidance says this can be a sign of a wrong answer key or unclear
        wording \citep{moodle_stats}.
  \item \textbf{Think about the pass mark.} If the goal is mastery, a pass mark
        of 4/10 may be too low, because most students stop there (H6). A
        higher threshold, or a small reward for reaching 8 or 10, could change
        this.
  \item \textbf{Use cheap warning signs.} A blank answer, or a first attempt
        under one minute, are easy to detect and are strongly linked to low
        performance (H5, H10). They could trigger a friendly message or extra
        support.
  \item \textbf{Do not judge students on Quiz 1 alone} (H8).
\end{enumerate}

\subsection{Reflection on our visualization choices}

Table~\ref{tab:vischoices} summarises the main chart used for each hypothesis.
A few lessons came out of the process.

\begin{itemize}
  \item \textbf{Our first idea was not always the right one.} For H1 we first
        drew a scatter plot of time against grade. With 30,000 points and only
        six grades it became six solid stripes and showed nothing. Replacing it
        with medians and IQR bars, split by attempt type, showed the real
        pattern, including the reversal.
  \item \textbf{Splitting the data was more important than the chart type.}
        Separating first attempts from re-attempts, and showing the quizzes
        separately, changed the answers to H1 and H5.
  \item \textbf{Highlighting one thing can reveal a story.} In H7, colouring
        only the full-mark re-attempts and greying the rest made the cluster of
        instant perfect re-attempts, and the sudden two-hour floor, visible at a
        glance. We would not have found the change in the quiz settings from a
        table.
  \item \textbf{We avoided some common mistakes on purpose:} no pie charts, no
        dual y-axes, no rainbow colour scales, no 3D. Where we needed both a
        trend and the number of cases (H5), we used two aligned panels.
\end{itemize}

\begin{table}[H]
\centering
\small
\caption{Summary of the main visualization for each hypothesis.}
\label{tab:vischoices}
\begin{tabularx}{\linewidth}{@{}l >{\raggedright\arraybackslash}p{2.6cm} >{\raggedright\arraybackslash}p{3.5cm} X@{}}
\toprule
\textbf{H} & \textbf{Chart} & \textbf{Main marks \& channels} & \textbf{Why this and not something else} \\
\midrule
H1 & Median + IQR dot-line, 2 panels & Points/lines; y-position (log time), x-position (grade), hue (quiz) & A scatter plot of 30k points on 6 grade values overplots; medians show the trend and IQR shows spread. \\
H2 & Heatmaps + tier dot plot & Cells with lightness + labels; points on a common x-scale & Heatmap gives an overview of 15 or 105 values; dot plot gives accurate tier comparison. \\
H3 & Line + CI, stacked bars, dot + CI & Lines (position), diverging hue for up/down, shape for tier method & Attempt number is ordered; diverging colours separate good and bad changes; two shapes show the extra check. \\
H4 & Labelled scatter; dot + bootstrap CI & Points: x = facility, y = time ratio, hue = quiz & Two numeric values for 15 items; medians with CIs, not bars, for skewed time. \\
H5 & Line + CI band, aligned count panel & Lines; x-position (log time bins), y-position (mean grade) & Shape of a trend needs a line; aligned panel instead of a dual axis. \\
H6 & Grouped bars; 100\% stacked bars & Length (proportion), ordered lightness & Proportions start at zero, so bars are correct; stacked bars compare two distributions better than pies. \\
H7 & Scatter over time (emphasis colour); stacked bars & Points: x = date, y = log wait; orange vs.\ grey & Shows each re-attempt, so a sudden rule change is visible. \\
H8 & Alluvial diagram + bar with CI & Blocks and ribbons: width = number of students & Designed to show movement between groups over steps \citep{rosvall2010mapping}. \\
H9 & Line + CI band, 2 panels & Lines; x-position (half-year), hue (quiz) & Ordered time; direct labels; two panels for two scales. \\
H10 & Box plots, dot + CI, bars & Hollow vs.\ filled for answered vs.\ blank & Same contrast repeated in three aligned panels. \\
\bottomrule
\end{tabularx}
\end{table}

\subsection{Limitations}

\begin{itemize}
  \item \textbf{No per-question time.} We only have the time for the whole
        attempt, so H4 could only be tested indirectly.
  \item \textbf{Time includes idle time.} Even under 60 minutes, some of the
        recorded time may be spent on other things. Our 60-minute cut-off is a
        judgement call, so we re-ran H1 with cut-offs of 30, 45, 90 and 120
        minutes. The correlations changed by at most 0.01, so the conclusions
        do not depend on this choice.
  \item \textbf{Inferred facts.} The pass mark of 4/10 and the two-hour delay
        are inferred from the data, not from the course settings. The
        patterns are very sharp, but they should be confirmed with the course
        team.
  \item \textbf{Unknown question content.} We do not know what each question
        asks or what type it is (for example multiple choice or typed answer),
        which might explain why only Q1.4 and Q2.5 are often left blank, and why
        some questions changed in 2024.
  \item \textbf{Correlation, not cause.} Apart from the before/after comparison
        in H7, all results are observational. For example, we cannot say that
        spending more time \emph{causes} a higher score.
  \item \textbf{Tier definition.} Tiers are based on first-attempt grades. We
        checked H3 with a second definition, but other definitions (for example
        best grade) could give slightly different groups.
  \item \textbf{Time zone and end of data.} The time zone of the timestamps is
        not given, so we did not analyse time of day. The last cohort (early
        2025) had less time to make re-attempts before the data was exported.
\end{itemize}

% ================================================================== 7
\section{Conclusion}
\label{sec:conclusion}

We cleaned and explored three years of attempt logs from three quizzes and used
visualizations to check ten hypotheses. Table~\ref{tab:summary} gives the
final answer for each one.

\begin{table}[H]
\centering
\small
\caption{Summary of all hypotheses.}
\label{tab:summary}
\begin{tabularx}{\linewidth}{@{}l X l >{\raggedright\arraybackslash}p{5.2cm}@{}}
\toprule
& \textbf{Hypothesis} & \textbf{Decision} & \textbf{Key evidence} \\
\midrule
H1 & Longer time $\rightarrow$ higher score & \partly & First attempts: weak $\rho$ = 0.08--0.23; re-attempts: $\rho$ = $-0.17$ to $-0.29$ \\
H2 & Some questions consistently harder & \accepted & Facility 9\%--80\%; same hardest question in 20 of 21 quiz half-years; $W$ = 0.85--0.95 \\
H3 & High performers improve steadily, low performers erratic & \rejected & 70--76\% of re-attempts go up in every tier; 18--20\% erratic in every tier \\
H4 & Hard questions take longer; strong students faster & \rejected & Time ratio rises with facility ($\rho$ = +0.72); Low tier fastest (indirect test) \\
H5 & Optimal time range & \partly & Under 1 min: $-1.6$ marks; small drop for very long attempts \\
\midrule
H6 & Re-attempt to pass, not to master & \accepted & Retry rate 98.7\% below 4/10 vs.\ 15.6\% at 4--8 \\
H7 & Instant re-attempts reuse feedback & \accepted & Gain +6.2 (70\% full marks, 1 min) before delay vs.\ +2.5 after \\
H8 & Good in Quiz 1 $\rightarrow$ good later & \partly & 84\% vs.\ 66\% reach 6--10 in Quiz 3; but $\rho \approx 0.2$ \\
H9 & Later cohorts score higher & \partly & Quiz 2: 3.6 $\rightarrow$ 6.1; Quiz 1 small rise; Quiz 3 no trend \\
H10 & Blank answer = low effort & \accepted & Faster (4.3 vs.\ 7.3 and 3.8 vs.\ 8.9 min); lower marks on other questions; mostly Low tier \\
\bottomrule
\end{tabularx}
\end{table}

Of the five given hypotheses, only H2 was fully accepted. The others were true
only in part or not at all once we separated first attempts from re-attempts
and looked at the right subsets. Our own hypotheses showed that the quiz rules,
especially the pass mark and the waiting time between attempts, explain a large
part of how students behave. Throughout the project, the most useful step was
not a fancy chart but understanding the data first and splitting it in the
right way; good visualization design then made the patterns easy to see and to
explain.

\textbf{Future work.} With per-question time logs we could test H4 directly.
With the course settings we could confirm the pass mark and the date of the
delay change. It would also be useful to link quiz behaviour with final exam
results, to see whether students who copied answers in 2022 learned less.

% ================================================================== references
\bibliographystyle{plainnat}
\bibliography{references}

% ================================================================== appendix
\appendix
\section{How to run the code}
\label{app:code}

The submitted zip file has the following structure:

\begin{verbatim}
quiz-data-visualization/
  README.md            how to install and run
  requirements.txt     Python packages
  run_all.py           runs everything (about 10 seconds)
  dataset/quiz1..3/    original CSV files (not changed)
  src/
    config.py          paths, cleaning thresholds, colours
    data_prep.py       loading, cleaning, attempt numbers, tiers
    stats_utils.py     confidence intervals and tests
    style.py           shared figure style
    eda.py             dataset overview (Figure 1)
    task1.py           H1-H5 (Figures 2-6)
    task2.py           H6-H10 (Figures 7-11)
  outputs/             created by run_all.py (figures, tables, results.json)
  report/              LaTeX source of this report
\end{verbatim}

To reproduce all results:
\begin{verbatim}
pip install -r requirements.txt
python run_all.py
\end{verbatim}
